\uuid{1pBQ}
\titre{Changement de variables}
\auteur{hanani ; mzari}
\organisation{exo7}
\licence{CC BY-NC-SA 4.0}
\datecreate{2020-01-20}
\niveau{L1}
\module{Analyse}
\chapitre{Calcul d'intégrales}
\sousChapitre{Primitives diverses}
\difficulte{3}
\theme{Primitives, changement de variables}
\exo7qcmid{lille-2.primitives.041}

\contenu{
\begin{question}
\qclassification{127.11}
\qtags{motcle=changement de variables}
Parmi les assertions suivantes, cocher celles qui sont vraies :
\begin{answers}
\bad{Le changement de variable $x=\sin t$ pour $t\in [-\pi/2,\pi/2]$ donne :
$$\int \sqrt{1-x^2}\mathrm{d}x=\int \sqrt{1-\sin ^2t}\, \mathrm{d}t=\int \cos t\, \mathrm{d}t.$$}
\good{Le changement de variable $x=\sin t$ pour $t\in [-\pi/2,\pi/2]$ donne :
$$\int \sqrt{1-x^2}\, \mathrm{d}x=\int \cos ^2t\, \mathrm{d}t.$$}
\bad{$\displaystyle \int \cos ^2t\, \mathrm{d}t=\frac{t}{2}-\frac{\sin (2t)}{4}+k=\frac{t}{2}-\frac{\sin t\cos t}{2}+k$, $k\in\Rr$.}
\good{Une primitive de $\displaystyle \sqrt{1-x^2}$ sur $[-1,1]$ est $\displaystyle \frac{\arcsin x}{2}+\frac{x\sqrt{1-x^2}}{2}$.}
\end{answers}
\begin{explanations}
Avec $x=\sin t$, on a : $\mathrm{d}x=\cos t\, \mathrm{d}t$ et $\sqrt{1-x^2}=\cos t$ car $t\in [-\pi/2,\pi/2]$. D'où
$$\int \sqrt{1-x^2}\, \mathrm{d}x=\int \cos ^2t\, \mathrm{d}t=\int \left(\frac{1}{2}+\frac{\cos (2t)}{2}\right)\mathrm{d}t=\frac{t}{2}+\frac{\sin (2t)}{4}+k,\; k\in \Rr.$$
Ainsi
$$\int \sqrt{1-x^2}\, \mathrm{d}x=\frac{t}{2}+\frac{\sin t\cos t}{2}+k=\frac{\arcsin x}{2}+\frac{x\sqrt{1-x^2}}{2}+k,\; k\in \Rr.$$
\end{explanations}
\end{question}
}
